\documentclass[a4paper]{article}
\usepackage[pdftex]{graphicx}
\usepackage[utf8]{inputenc}
\usepackage{enumerate}
\usepackage{amssymb}
\usepackage{icomma}
\usepackage{href-ul}
\hypersetup{
	colorlinks=true,
	linkcolor=blue,
	urlcolor=blue}
\usepackage{siunitx}
\sisetup{locale=DE} 
\usepackage{geometry}
\geometry{a4paper, top=15mm, left=15mm, right=15mm, bottom=15mm,
	headsep=10mm, footskip=12mm}

\begin{document}
{\bf Potenzielle Energie, Hubarbeit, Leistung, Wirkungsgrad}

\begin{enumerate}[1.]
{\bf \item Hubarbeit, Leistung: Fülle die fehlenden Werte der Tabelle aus!}

\renewcommand{\arraystretch}{2}
\begin{tabular}{|p{1,7 cm}|p{1,7 cm}|p{1,7 cm}|p{1,7 cm}|p{1,7 cm}|p{1,7 cm}|p{1,7 cm}|p{1,7 cm}|}
\hline
             & Aufzug	& Berg\-wanderer & E-Bike & Donau\-staustufe & Hau den Lukas & Silvester\-rakete & Ariane 5  \\
\hline
Masse $m$    &$\SI{450}{\kilogram}$ &$\SI{75}{\kilogram}$ & &$\SI{340}{\tonne}$ &$\SI{0,25}{\kilogram}$ &$\SI{80}{\gram}$ & \\
\hline
Höhe $h$     &$\SI{40}{\meter}$ &  &$\SI{220}{\meter}$ &  &$\SI{3,5}{\meter}$ & &$\SI{500}{\kilo\meter}$ \\
\hline
$W_{hub}$    &  &$\SI{700}{\kilo\joule}$ & &$\SI{25}{\mega\joule}$ & &$\SI{39}{\joule}$ &$\SI{3,8}{\tera\joule}$ \\
\hline
Zeit $t$     &$\SI{30}{\second}$ &  &$\SI{40}{\minute}$ &$\SI{1,0}{\second}$ & &  &$\SI{10}{\minute}$ \\
\hline
Leistung $P$ &  &$\SI{56}{\watt}$ &$\SI{94}{\watt}$ & &$\SI{10}{\watt}$ &$\SI{7,8}{\watt}$ & \\
\hline
\end{tabular}

{\bf \item Wirkungsgrad: Die folgenden Aufgaben beziehen sich auf die Beispiele der Tabelle. }

\begin{enumerate}[a)]
	\item Welche elektrische Leistung nimmt der Aufzugsmotor auf, wenn sein Wirkungsgrad 90\% beträgt?
	\item Wieviel KiloJoule an Nahrung muss der Bergwanderer für seine Tour essen, wenn der Wirkungsgrad seiner Muskeln 30\% beträgt?
	\item Beim Hau den Lukas ist die Schlagenergie des Hammers etwa $\SI{80}{\joule}$. Wieviel Prozent davon werden in potenzielle Energie des kleinen Gewichtes umgewandelt?
	\item Die Leistung des fließenden Wassers in der Donaustaustufe wird zu 80\% in elektrische Leistung umgewandelt. Wie groß ist diese nutzbare Leistung dann ?

\end{enumerate} 


\end{enumerate} 

\fbox{
	\begin{minipage}{0.5\textwidth}
		Zur Lösung bitte \href{https://www.okuyakl.de/phys/p8mghL077/ll077.pdf}{hier klicken} oder den QR-Code scannen.\\
	Weitere Arbeitsblätter gibt es unter 
	
	\href{https://www.okuyakl.de}{www.okuyakl.de}
	\end{minipage}
	\hfill
	\begin{minipage}{0.4\textwidth}
		\includegraphics[width=1.5 cm]{../../viecher/zwe03}
		\includegraphics[width=3 cm]{qr077}
		\includegraphics[width=2 cm]{../../viecher/afanticon1}
		
\end{minipage}}
\end{document}

